\scanpage{182}

isomorphe à une structure produit,
\[
X_i\times Z
\]
(avec les $X_j$ correspondant à des $Z_j\subset Z$, $\overline{X_j}$ à $\overline{Z_j}$, $X_i$ à un $z_0\in Z$ contenu dans les $\overline{Z_j}$), et les $\pi_0(Z_j;\overline{Z_j},z_0)$ sont de cardinal 1, i.e. il existe \deleted{\unsure{touj}} un syst. fond. de voisinages de $z_0$ dans $Z$ dont l'intersection avec les $Z_j$ est connexe)

\margin{\uline{NB} les $Z_j$ loc. connexes}

\uline{Remarque:} Si $U$ est un ouvert de $X$, alors l'ens des composantes connexes des $U\cap X_i$ forme \inserted{encore} une stratification \inserted{globale} de $U$. Donc \deleted{l'application} \inserted{pour un} ouvert $U$ variable de $X$,
\[
U\unsure{\mapsto}\mathrm{Strat}_{\mathrm{glob}}(U)
\]
forme un préfaisceau sur $X$. Les exemples de d) montrent que ce préfaisceau n'est pas \deleted{\unsure{séparé}} \inserted{un faisceau} (contre exemple encore plus simple: les feuilletages!)

\margin{j'ai l'impression cependant qu'une stratification globale \unsure{est connue quand on la connaît localement}.}

Les sections du faisceau associé s'appellent les \inserted{\uline{structures de}} \uline{stratification locale} sur $X$. \unsure{Comparer} avec l'exercice de SGA (théorie des Topos) sur les relations d'équivalence locales --- les stratifications en sont des exemples très

\margin{181}
